\documentclass[12pt,a4paper]{article}

\usepackage[utf8]{inputenc}
\usepackage[T1]{fontenc}
\usepackage[ngerman]{babel}
\usepackage{amsmath}
\usepackage{amssymb}
\usepackage{amsfonts}
\usepackage{mathtools}
\usepackage[left=3cm,right=2.5cm,top=2.5cm,bottom=2.5cm]{geometry} 
\usepackage{setspace}          
\onehalfspacing                
\usepackage{microtype}
\usepackage[autostyle=true,german=quotes]{csquotes}         

\usepackage[document]{ragged2e}
\hyphenpenalty=10000
\exhyphenpenalty=10000
\sloppy

\usepackage{graphicx}
\usepackage{booktabs}
\usepackage{caption}
\usepackage{listings}

\lstset{
    basicstyle=\ttfamily\small,
    breaklines=true,
    showstringspaces=false,
    language=C
}

\title{Systemprogrammierung Prüfung, WS25/26}
\author{Mathias Lampert}
\date{\today}

\begin{document}

\maketitle
\tableofcontents
\newpage

\section{C-Sprache}

\subsection{Bytespeicher für uint8\_t* und uint16\_t*}
Beide Pointer belegen den gleichen Speicherplatz. Die Größe eines Pointers hängt von der Systemarchitektur ab. Ein 32-Bit-System verwendet 4 Bytes für eine Adresse. Ein 64-Bit-System verwendet 8 Bytes. Ein Pointer speichert ausschließlich die Speicheradresse. Der Datentyp des Ziels definiert die Interpretation der Daten. Der Datentyp ändert die Größe des Pointers nicht.

\subsection{Einsatz von static in C}
\begin{itemize}
    \item Lokale Variablen: Die Variable behält ihren Wert über das Funktionsende hinaus. Das System speichert sie im Datensegment statt auf dem Stack. Beispiel: \texttt{static int counter = 0;} in einer Funktion.
    \item Globale Variablen: Das System beschränkt den Zugriff auf die aktuelle Quelldatei. Andere Dateien sehen diese Variable nicht. Beispiel: \texttt{static int file\_internal\_status;}
    \item Funktionen: Die Funktion ist nur innerhalb der eigenen C-Datei aufrufbar. Das verhindert Namenskonflikte beim Linken. Beispiel: \texttt{static void helper(void);}
\end{itemize}

\subsection{Endianness}
Endianness definiert die Byte-Reihenfolge im Speicher. Little Endian speichert das niederwertigste Byte an der kleinsten Adresse. RISC-V nutzt dieses Format. Big Endian speichert das höchstwertigste Byte zuerst. Portierungsprobleme treten auf; wenn Software Daten über Netzwerke schickt oder binäre Dateiformate zwischen Systemen austauscht. Sie müssen die Byte-Reihenfolge in diesen Fällen manuell umkehren.

\newpage
\section{Konfiguration}
\begin{lstlisting}
#include <stdio.h>
#include <stdint.h>

typedef struct {
    uint8_t mode;
    uint8_t level;
    uint8_t threshold;
    uint8_t flags;
} SensorConfig;

typedef union {
    uint32_t raw;
    SensorConfig config;
} SensorData;

void printSensorData(SensorData *data) {
    printf("Raw-Wert: 0x%08X\n", data->raw);
    printf("Mode: %u, Level: %u, Threshold: %u, Flags: %u\n", 
            data->config.mode, 
            data->config.level, 
            data->config.threshold, 
            data->config.flags);
}
\end{lstlisting}

\newpage
\section{Bitmaskierung}

\subsection{Aufgabe: setCtrl-Funktion}
\begin{lstlisting}
void setCtrl(uint32_t* pCtrl, bool enableMotor, bool enableLight, bool enableSensors) {
    uint32_t mask = (1 << 0) | (1 << 8) | (1 << 16);
    *pCtrl &= ~mask; 
    if (enableMotor)   *pCtrl |= (1 << 0);
    if (enableLight)   *pCtrl |= (1 << 8);
    if (enableSensors) *pCtrl |= (1 << 16);
}
\end{lstlisting}

\subsection{Aufgabe: Rotate Right}
\begin{lstlisting}
uint16_t rotateRight(uint16_t value, uint8_t n) {
    n %= 16;
    return (value >> n) | (value << (16 - n));
}
\end{lstlisting}

\newpage
\section{Dynamische Speicherverwaltung}

\subsection{Aufgabe: mem\_calcStatistics}
\begin{lstlisting}
void mem_calcStatistics(float* avgAlloc, float* avgFree) {
    uint32_t sumAlloc = 0, countAlloc = 0;
    uint32_t sumFree = 0, countFree = 0;
    uint32_t offset = 0;

    while (offset < HEAPSIZE) {
        struct Header* h = (struct Header*)&gHeap[offset];
        if (h->flags & 0x01) {
            sumAlloc += h->length;
            countAlloc++;
        } else {
            sumFree += h->length;
            countFree++;
        }
        offset += sizeof(struct Header) + h->length;
    }
    *avgAlloc = countAlloc ? (float)sumAlloc / countAlloc : 0;
    *avgFree = countFree ? (float)sumFree / countFree : 0;
}
\end{lstlisting}

\newpage
\section{Multitasking}

\subsection{Semaphor}
Ein Semaphor steuert den Zugriff auf Ressourcen. Er nutzt einen Zähler.
\begin{itemize}
    \item P() (Take): Verringert den Zähler. Ist der Zähler Null; wartet der Task.
    \item V() (Give): Erhöht den Zähler. Das System weckt einen wartenden Task auf.
\end{itemize}

\subsection{Semaphor in FreeRTOS}
\begin{itemize}
    \item xSemaphoreCreateCounting: Erstellt einen zählenden Semaphor. uxMaxCount legt den maximalen Zählerstand fest. uxInitialCount definiert den Startwert.
    \item xSemaphorGive: Gibt den Semaphor frei; erhöht den Zähler. Rückgabe ist pdTRUE bei Erfolg.
    \item xSemaphoreTake: Reserviert den Semaphor; verringert den Zähler. xTicksToWait definiert die maximale Wartezeit.
\end{itemize}
\newpage
\subsection{Producer/Consumer Implementierung}
\begin{lstlisting}
SemaphoreHandle_t xSem;

void producerTaskFunc(void * pvParameters) {
    for (;;) {
        xSemaphoreGive(xSem);
        vTaskDelay(pdMS_TO_TICKS(30));
        xSemaphoreGive(xSem);
        vTaskDelay(pdMS_TO_TICKS(50));
    }
}

void consumerTaskFunc(void * pvParameters) {
    for (;;) {
        if (xSemaphoreTake(xSem, portMAX_DELAY)) {
            printf("Event verarbeitet\n");
        }
    }
}

void setupAndStartTasks(void) {
    xSem = xSemaphoreCreateCounting(10, 0);
    xTaskCreate(producerTaskFunc, "Prod", 2048, NULL, 2, NULL);
    xTaskCreate(consumerTaskFunc, "Cons", 2048, NULL, 1, NULL);
}
\end{lstlisting}

\subsection{Scheduling Diagramm}
Ein Tick entspricht 10 ms. Der Idle-Task nutzt die Zeit; wenn Producer und Consumer auf ihre Delays oder Semaphore warten.
\begin{itemize}
    \item 0-10ms: Producer läuft; Consumer läuft; beide blockieren. Idle übernimmt.
    \item 30ms: Producer wacht auf; gibt Semaphor; blockiert. Consumer wacht auf; arbeitet; blockiert.
\end{itemize}

\newpage
\subsection{Producer/Consumer Erweiterung}
Sie können folgende Strukturen nutzen:
\begin{itemize}
    \item Message Queues für den Datenaustausch.
    \item Stream Buffer für kontinuierliche Daten.
    \item Task Notifications für einfache Signale.
\end{itemize}

\end{document}